% Referencias verificadas el 10 de septiembre de 2026. No modifica datos.
\begin{thebibliography}{99}
\setlength{\itemsep}{0.10\baselineskip}
\phantomsection
\addcontentsline{toc}{section}{Références}

\bibitem{mustovicary}
B. Musto et J. Vicary,
``Quantum Latin squares and unitary error bases'',
\emph{Quantum Information and Computation} \textbf{16},
numéros 15--16, 1318--1332 (2016).
\href{https://doi.org/10.26421/QIC16.15-16-4}{doi:10.26421/QIC16.15-16-4}.

\bibitem{delascuevas2020}
G. De las Cuevas, T. Drescher et T. Netzer,
``Quantum magic squares: Dilations and their limitations'',
\emph{Journal of Mathematical Physics} \textbf{61}, 111704 (2020).
\href{https://doi.org/10.1063/5.0022344}{doi:10.1063/5.0022344}.
\href{https://arxiv.org/abs/1912.07332}{Prépublication des auteurs : arXiv:1912.07332}.

\bibitem{delascuevas2023}
G. De las Cuevas, T. Netzer et I. Valentiner-Branth,
``Magic squares: Latin, semiclassical, and quantum'',
\emph{Journal of Mathematical Physics} \textbf{64}, 022201 (2023).
\href{https://doi.org/10.1063/5.0127393}{doi:10.1063/5.0127393}.
\href{https://arxiv.org/abs/2209.10230}{Prépublication des auteurs : arXiv:2209.10230}.

\bibitem{codata2018}
E. Tiesinga, P. J. Mohr, D. B. Newell et B. N. Taylor,
``CODATA recommended values of the fundamental physical constants: 2018'',
\emph{Reviews of Modern Physics} \textbf{93}, 025010 (2021).
\href{https://doi.org/10.1103/RevModPhys.93.025010}{doi:10.1103/RevModPhys.93.025010}.
\href{https://physics.nist.gov/cuu/Constants/archive2018.html}{NIST : archive de l’ajustement de 2018};
\href{https://physics.nist.gov/cuu/Constants/ArchiveASCII/allascii_2018.txt}{tableau ASCII archivé}.

\bibitem{codata2022}
P. J. Mohr, D. B. Newell, B. N. Taylor et E. Tiesinga,
``CODATA recommended values of the fundamental physical constants: 2022'',
\emph{Reviews of Modern Physics} \textbf{97}, 025002 (2025).
\href{https://doi.org/10.1103/RevModPhys.97.025002}{doi:10.1103/RevModPhys.97.025002}.
\href{https://physics.nist.gov/cuu/Constants/article2022.html}{NIST : documentation de l’ajustement de 2022}.
Les comparaisons de l’article conservent identifiés les ajustements de 2018 et 2022.

\bibitem{pdg2026}
F. Takahashi et al. (Particle Data Group),
``Review of Particle Physics'',
\emph{International Journal of Modern Physics A} \textbf{41}, 2630011 (2026).
\href{https://doi.org/10.1142/S0217751X26300115}{doi:10.1142/S0217751X26300115}.
\href{https://pdg.lbl.gov/2026/}{Édition officielle de 2026};
\href{https://pdg.lbl.gov/2026/api/index.html}{documentation d’accès programmatique}.
Le catalogue de cet article conserve l’instantané identifié dans les sources
comme PDG 2026, avec un état documentaire arrêté au 15 janvier 2026 et une consultation consignée
le 26 juillet 2026. Ses 471 enregistrements de masse et 384 de largeur, ainsi que
l’empreinte de provenance, sont maintenus dans le supplément de données ; la référence
web ne remplace pas cet instantané.

\end{thebibliography}
